\section{Internal group structure}
\label{sec:group-structure}

Write
\[
A=\operatorname{Aut}(G60).
\]
By Section~\ref{sec:complete-enumeration}, the group $A$ is now known
completely as a set of $480$ native permutations. This permits exact
calculation of its center, derived series, element orders, and
index-two subgroups.

These invariants do more than distinguish candidate groups of order
$480$. They expose the way the central deck involution is embedded in
the full symmetry group.

\subsection{The center}

The center is computed directly by testing commutation with every group
element:
\[
Z(A)
=
\{z\in A:zg=gz\text{ for every }g\in A\}.
\]

The calculation gives
\[
|Z(A)|=2.
\]

Since the canonical deck involution $a$ is nontrivial and central by
Proposition~\ref{prop:deck-central}, it follows that
\[
Z(A)=\langle a\rangle\cong C_2.
\]

\begin{proposition}
\label{prop:center}
The center of the native automorphism group is
\[
Z(A)\cong C_2,
\]
generated by the canonical deck involution.
\end{proposition}

Thus the deck symmetry is the unique nonidentity central automorphism of
$G60$.

\subsection{The derived series}

Let
\[
A'=[A,A]
\]
denote the derived subgroup. Exact commutator generation gives
\[
|A'|=120.
\]

The subgroup $A'$ has center of order $2$, and its quotient by that
center is isomorphic to $A_5$. More precisely,
\[
A'\cong A_5\times C_2.
\]

The second derived subgroup satisfies
\[
A''=[A',A']\cong A_5.
\]

\begin{theorem}
\label{thm:derived-series}
The derived series begins
\[
A
\trianglerighteq
A_5\times C_2
\trianglerighteq
A_5
\trianglerighteq
A_5.
\]
In particular,
\[
A'\cong A_5\times C_2
\qquad\text{and}\qquad
A''\cong A_5.
\]
\end{theorem}

The deck involution belongs to the first derived subgroup:
\[
a\in A'.
\]
This point is structurally important. Although $a$ is central in the
full group, it is not detected as an independent factor in the
abelianization.

\subsection{The abelianization}

Since
\[
|A|=480
\qquad\text{and}\qquad
|A'|=120,
\]
the abelianization has order
\[
|A/A'|=4.
\]

Direct computation of the quotient shows that it is not cyclic:
\[
A/A'\cong C_2\times C_2.
\]

\begin{corollary}
\label{cor:abelianization}
The abelianization of the native automorphism group is
\[
A^{\mathrm{ab}}
=
A/A'
\cong C_2\times C_2.
\]
\end{corollary}

Accordingly, $A$ has three index-two subgroups, corresponding to the
three nontrivial characters
\[
A\longrightarrow C_2.
\]

\subsection{Element-order profile}

The complete element-order distribution is
\[
\begin{array}{c|rrrrrrrr}
\text{order}
&1&2&3&4&5&6&10&12\\
\hline
\text{count}
&1&83&20&140&24&100&72&40.
\end{array}
\]

The counts sum to
\[
1+83+20+140+24+100+72+40=480.
\]

In particular, $A$ contains no elements of order $8$, $15$, $20$, or
$24$. The simultaneous presence of elements of orders $4$ and $12$,
together with the derived-series data, sharply constrains possible
extensions of $S_5$-type symmetry by a group of order $4$.

\subsection{The three index-two subgroups}

Because
\[
A/A'\cong C_2\times C_2,
\]
there are exactly three subgroups of index two.

One is a split subgroup
\[
H_{\mathrm{split}}\cong S_5\times C_2.
\]

A second is the sign-pullback subgroup
\[
H_{\mathrm{twist}}
\cong
S_5\times_{\operatorname{sgn}} C_4,
\]
which may be written
\[
H_{\mathrm{twist}}
=
\left\{
(\sigma,u)\in S_5\times C_4:
\operatorname{sgn}(\sigma)=u\bmod 2
\right\}.
\]
This subgroup is nonsplit over the relevant $S_5$ quotient.

The third index-two subgroup has center of order $4$ and element-order
profile
\[
\begin{array}{c|rrrrrrr}
\text{order}
&1&2&3&5&6&10\\
\hline
\text{count}
&1&63&20&24&60&72.
\end{array}
\]

These three subgroups distinguish the native group from simpler direct
products and from an extension determined only by its order.

\subsection{Structural summary}

The independently established invariants are
\[
|A|=480,
\]
\[
Z(A)\cong C_2,
\]
\[
A'\cong A_5\times C_2,
\]
\[
A''\cong A_5,
\]
and
\[
A/A'\cong C_2\times C_2.
\]

Together with the split and nonsplit index-two subgroups, these data
point to a fiber product of $S_5$ with a dihedral group of order $8$.
The next section constructs that model and determines which quotient
character of the dihedral factor gives the native group.
